% OWNER: sections/s3_design.tex. Edit only this file.
\seclabel
\claim{\secnum{3}{sec:design}\quad The wrapper keeps one login session alive; strategy and market data stay with the caller}

\lead{The wrapper has 5 goals and 5 non-goals, and the largest exclusion is market data (\figref{goals}).} The contest server reports a price only for the stocks the portfolio holds (\code{close} in each position record), and the club's market data comes from the Korea Investment \& Securities API. Folding that API into the wrapper would tie a contest client to a broker's authentication and rate limits, so the caller passes prices in and the wrapper stays a client of the contest site alone. A reader may ask how \code{plan()} then sizes an order in a stock the portfolio does not hold: the caller passes that price through \code{prices=}, and an order without a price is kept with a warning (Section~\ref{sec:oneshare}).

% OWNER: figs/goals.tex. Edit only this file.
\fig[wide]{goals}{5 goals and 5 non-goals bound the wrapper}[Scope of \code{pytimefolio}: each goal with the measure that shows it is met]{%
{\tabfont\begin{tabularx}{\linewidth}{>{\RaggedRight\arraybackslash}p{7mm} >{\RaggedRight\arraybackslash}p{70mm} >{\RaggedRight\arraybackslash}X}
\thd{ID} & \thd{Goal} & \thd{Measure of success}\\
\hl G1 & Stay logged in for the whole contest without a person & No manual login between the start and 30 Nov; every renewal and login logged\\\rowsep
G2 & Give the strategy the portfolio's current state & State no older than one poll (3\,s) and equal to a fresh snapshot after every delta; reads, orders and \code{wait_done()} refuse a state older than 26\,s (\code{StaleFeed})\\\rowsep
G3 & Place, stop and track orders within the organiser's limits and the contest's guidelines & At most 80 orders a day (new orders plus stopped orders); at least 1\,s between writes; an order whose reply is lost is never resent; \code{plan()} refuses targets that break the total, universe, stock or sector rule\\\rowsep
G4 & Polling never delays an order & Feed polling and orders share one event loop, and no order waits for a feed reply (Section~\ref{sec:live})\\\rowsep
G5 & Stop on ambiguity instead of guessing & Raises \code{AmbiguousState} when the portfolio rule matches 0 entries or more than 1, when \code{currBizDay} differs from the KST date, or when a placed order is found neither in the feed within 15\,s nor in REST \code{Orders} (the order is not resent); raises \code{RequestError} on a 400\\
\specialrule{0.8pt}{1pt}{1pt}
\thd{ID} & \thd{Non-goal} & \thd{Reason}\\
N1 & Market data (prices, order books) & The caller supplies prices; folding in another API adds another login and another set of limits\\\rowsep
N2 & Strategy and portfolio construction & The wrapper turns target weights into orders; choosing targets is the caller's\\\rowsep
N3 & Browser automation & Every action would wait on page rendering and break when the layout changes; the page's own HTTP requests take 16 to 39\ms{} (Section~\ref{sec:transport})\\\rowsep
N4 & Several accounts or other participants' portfolios & One account per process; other participants' data is not ours to query\\\rowsep
N5 & Probing rate limits & The organiser may block an IP address for bursts of traffic\\
\end{tabularx}}}{G = goal, N = non-goal. Shaded: G1 is shown only by running until 30 Nov; the live tests of Section~\ref{sec:live} cover one morning. The 80-order cap leaves a margin of 20 under the organiser's threshold of 100 orders a day.}{DART Quant Team}


\lead{Where the server or the organiser sets a limit, the wrapper sets its own parameter a margin inside it: 80 orders a day against the organiser's 100, renewal at 720\,s against a 900\,s token, \code{keepalive_expiry} 10\,s against an idle cut-off of 15 to 30\,s.} Appendix~\ref{app:decisions} lists each parameter with the measurement behind it and the condition that would make us change it.

\lead{A trading day runs in 5 phases, and all login and set-up work falls between the business-day rollover (06:02 to 07:07) and 08:50 (\figref{schedule}).} During the trading session (09:00 to 15:20) the wrapper only polls the live feed and sends the caller's orders. It has no renewal timer: the first REST request more than 720\,s after the token was issued, by the server's clock, renews it, and the live feed asks for a token only when it connects or reconnects.

% OWNER: figs/schedule.tex. Edit only this file.
\fig[wide]{schedule}{The wrapper finishes set-up before 08:50 and only polls and trades during the session}[A trading day for a strategy that rebalances once at the open]{%
{\tabfont\begin{tabularx}{\linewidth}{>{\RaggedRight\arraybackslash}p{20mm} >{\RaggedRight\arraybackslash}p{26mm} >{\RaggedRight\arraybackslash}X}
\thd{Time (KST)} & \thd{Phase} & \thd{Wrapper actions}\\
07:10 & Start & Renew the token from the stored cookie jar, or log in if renewal fails; choose the portfolio; orders are refused until \code{currBizDay} equals the KST date\\\rowsep
07:15 & Load & Open the live feed and take a snapshot; load \code{Universe} for each stock's sector; count the day's orders on the server\\\rowsep
08:00 to 08:50 & Plan (caller) & The caller computes target weights from its own prices; \code{plan()} checks them against the universe, the total, the stock and sector limits and one share's weight, and converts them into orders; \code{guidelines()} reports the turnover and small-cap totals\\\rowsep
08:50 & Submit & \code{rebalance()} re-counts the day's orders on the server, checks the daily cap of 80, then sends sells first, at least 1\,s apart; each order is confirmed when it appears in the feed's \code{ords}\\\rowsep
09:00 to 15:20 & Session & Poll the feed every 3\,s; the first REST call after 720\,s renews the token, and the feed asks for a token only when it reconnects; stop orders on the caller's request\\
\end{tabularx}}}{The caller chooses the times; apart from the 3\,s poll, the wrapper acts only when called. An order placed at 08:30 was \code{Working} by 08:33, began filling at 09:00, the session's start, and was \code{Done} at 09:08:07 (Section~\ref{sec:orders}).}{DART Quant Team}


\lead{We chose direct HTTP requests, the same ones the page sends, because a reused connection answers one from Python in 31 to 39\ms{} (Section~\ref{sec:transport}) and no page has to render (\figref{alternatives}).} Other contest tools we found on GitHub drive the web page with Playwright, which avoids reverse-engineering the requests, but every action waits for the page to render, the script breaks when the layout changes, and the server needs a browser installed.

% OWNER: figs/alternatives.tex. Edit only this file.
\fig[wide]{alternatives}{Direct HTTP requests from one \code{asyncio} thread; an unconfirmed order is never resent}[Alternatives considered for each design decision]{%
{\tabfont\begin{tabularx}{\linewidth}{>{\RaggedRight\arraybackslash}p{27mm} >{\RaggedRight\arraybackslash}p{52mm} >{\RaggedRight\arraybackslash}X}
\thd{Decision} & \thd{Chosen} & \thd{Alternatives and why they lost}\\
How to reach the server & HTTP requests as the page makes them: 16.0\ms{} per reused request with curl, 31 to 39\ms{} from Python & Browser automation (Playwright): every action waits on page rendering, breaks when the layout changes and needs a browser on the run host\\\rowsep
Concurrency & \code{asyncio} in one thread & Synchronous requests: an order would wait while a feed request waits for its reply. Threads: the feed thread can change the state between any 2 statements, so every read of the state needs a lock; \code{asyncio} switches tasks only at \code{await}. Both tasks spend their time waiting on the network, so the GIL costs nothing\\\rowsep
Feed transport & WebSocket only & Server-Sent Events and long polling, offered in the server's negotiation reply: long polling adds 1 HTTP request per message. Neither is implemented\\\rowsep
Token on the WebSocket & \code{Authorization} header & In the URL, as browsers must send it: the token would end up in access logs\\\rowsep
Lost \code{AddOrder} reply & Never resend: look for the order in the feed for 15\,s, then in REST \code{Orders}; if it is in neither, raise \code{AmbiguousState}. Resend once only after a connection failure that never reached the server & Resend at once: risks a duplicate order, since the reply carries no order ID\\
\end{tabularx}}}{``Chosen'' is what \code{pytimefolio} 0.1.0 implements. Request times are from Section~\ref{sec:transport}; we did not time browser automation. The negotiation reply of the live feed (Section~\ref{sec:feed}) lists WebSockets, Server-Sent Events and long polling.}{DART Quant Team}


\lead{All work runs as coroutines on 1 asyncio event loop in 1 thread, and 3 locks guard the only steps that must not interleave (\figref{asyncio}).} A coroutine gives up the thread only at an \code{await}, so a delta merge, which has none, never sees another coroutine's half-written state, and the feed task and the caller's coroutines share the state without a lock. The work is waiting on sockets: 1 request takes 31 to 39\ms{} of which the thread is busy for under 1\ms{}, so 1 thread is enough. The 3 locks cover the steps that do await in the middle.
\begin{itemize}
\item Renewal is single-flight. A coroutine that finds the token older than 720\,s takes the token lock and checks the age again inside it; a second coroutine that was waiting for the lock finds a new token and returns without renewing. In the offline test 5 concurrent \code{token()} calls sent 1 \code{SilentRefresh}.
\item Writes pass a gate. The gate lock lets 1 write out at a time, at least 1\,s after the previous one, and only while the day's count plus the request stays at or under 80.
\item Placing an order holds a lock across the write and the search for the new order. The placing coroutine looks for the order in the feed's state every 0.5\,s for up to 15\,s; the feed itself refreshes every 3\,s. Two identical orders placed together therefore cannot both bind the first new record's order ID.
\end{itemize}

% OWNER: figs/asyncio.tex. Edit only this file.
% Gantt block: \tfgb{y}{x0}{x1}{fill}{text colour}{label}
\providecommand{\tfgb}[6]{\fill[#4] (#2,#1-2) rectangle (#3,#1+2);
  \node[font=\fontsize{7}{8.4}\selectfont,text=#5,inner sep=0pt] at ({(#2+#3)/2},#1) {#6};}
\fig[wide]{asyncio}{One thread runs the coroutines in turn; the token lock lets only one of them renew}[Gantt chart of the event loop while \code{rebalance()} places 2 orders and \code{guidelines()} runs beside it]{%
\begin{tikzpicture}[x=1mm,y=1mm,
  f/.style={font=\fontsize{7}{8.4}\selectfont},
  lane/.style={f,anchor=east,align=right,text=navy},
  note/.style={f,text=greyD,inner sep=0.5pt},
  ptitle/.style={font=\fontsize{7.5}{9}\selectfont\bfseries,text=navy,anchor=west},
  ev/.style={-{Stealth[length=1.4mm]},line width=0.5pt,navyM,dashed}]
% ====================================================== panel (a): 0 to 145 ms, 1 mm per ms
\node[ptitle] at (-30,8) {(a) First 145\,ms: one coroutine runs at a time; G waits for the token lock instead of renewing again};
\node[lane] at (-2,0)   {\textbf{Event loop}\\{\color{greyD}one thread}};
\node[lane] at (-2,-9)  {\textbf{Feed task}\\{\color{greyD}\code{hub.stream()}}};
\node[lane] at (-2,-21) {\textbf{Coroutine R}\\{\color{greyD}\code{rebalance()}}};
\node[lane] at (-2,-32) {\textbf{Coroutine G}\\{\color{greyD}\code{guidelines()}}};
\node[lane] at (-2,-42) {\textbf{Token lock}\\{\color{greyD}\code{Auth._lock}}};
\node[lane] at (-2,-51) {\textbf{Place lock}\\{\color{greyD}\code{Orders._place_lock}}};
\begin{scope}
% event loop strip
\draw[greyM,line width=0.4pt] (0,-2) rectangle (145,2);
\tfgb{0}{0}{3}{navy}{white}{R}     \tfgb{0}{3}{6}{navy}{white}{G}
\tfgb{0}{35}{38}{navy}{white}{R}   \tfgb{0}{38}{41}{navy}{white}{G}
\tfgb{0}{70}{73}{navy}{white}{R}   \tfgb{0}{73}{76}{navy}{white}{G}
\tfgb{0}{105}{108}{navy}{white}{R} \tfgb{0}{140}{143}{navy}{white}{R}
\foreach \a/\b in {6/35,41/70,76/105,108/140} {\node[note] at ({(\a+\b)/2},0) {idle};}
% feed task
\tfgb{-9}{0}{145}{navyL}{navy}{\code{sleep(3)} between feed requests; next request in panel (b)}
% coroutine R
\tfgb{-21}{0}{3}{navy}{white}{}
\tfgb{-21}{3}{35}{navyM}{white}{\code{SilentRefresh}}
\tfgb{-21}{35}{38}{navy}{white}{}
\tfgb{-21}{38}{70}{navyM}{white}{\code{Orders} (re-count)}
\tfgb{-21}{70}{73}{navy}{white}{}
\tfgb{-21}{73}{105}{navyM}{white}{\code{ValidateProduct}}
\tfgb{-21}{105}{108}{navy}{white}{}
\tfgb{-21}{108}{140}{navyM}{white}{\code{AddOrder}}
\tfgb{-21}{140}{143}{navy}{white}{}
\fill[navyL] (143,-23) rectangle (145,-19);
\draw[decorate,decoration={brace,amplitude=1.2mm},navy,line width=0.5pt] (0,-18.3) -- (36.5,-18.3);
\node[f,text=navy,anchor=south west] at (0,-17.3) {token older than 720\,s: lock, re-check, renew};
\draw[decorate,decoration={brace,amplitude=1.2mm},navy,line width=0.5pt] (105,-18.3) -- (108,-18.3);
\node[f,text=navy,anchor=south east] at (145,-17.3) {write gate: at least 1\,s since the last write, under 80 orders today};
% coroutine G
\tfgb{-32}{3}{6}{navy}{white}{}
\tfgb{-32}{6}{38}{greyM}{navy}{waits for the lock}
\tfgb{-32}{38}{41}{navy}{white}{}
\tfgb{-32}{41}{73}{navyM}{white}{\code{Violations}}
\tfgb{-32}{73}{76}{navy}{white}{}
\node[note,anchor=north west] at (3,-34.3) {token due, lock busy};
\node[note,anchor=north west] at (38,-34.3) {re-check inside the lock: token already new, no 2nd renewal};
\node[note,anchor=west] at (77,-32) {returns};
% token lock
\tfgb{-42}{0}{36.5}{navy}{white}{held by R}
\node[note,anchor=west] at (38,-42) {free; G enters and leaves at once};
% place lock
\tfgb{-51}{70}{145}{navy}{white}{held by R from the checks through write and find}
% axis
\draw[greyD,line width=0.4pt] (0,-54.6) -- (145,-54.6);
\foreach \t in {0,35,70,105,140} {\draw[greyD,line width=0.4pt] (\t,-54.6) -- (\t,-55.6); \node[f,text=greyD,anchor=north] at (\t,-55.6) {\t};}
\node[f,text=greyD,anchor=north west] at (0,-58.8) {ms after R and G start together; each request drawn 32\,ms (measured 31 to 39\,ms); dark slices drawn 3\,ms wide};
\end{scope}
% ====================================================== panel (b): 0 to 10 s, 12 mm per s
\begin{scope}[yshift=-77mm]
\node[ptitle] at (-30,9) {(b) First 10\,s: each order is confirmed by finding it in the feed's state; the write gate spaces the 2 orders};
\def\s{12}
\node[lane] at (-2,0)   {\textbf{Feed task}\\{\color{greyD}\code{hub.stream()}}};
\node[lane] at (-2,-12) {\textbf{Coroutine R}\\{\color{greyD}\code{rebalance()}}};
\node[lane] at (-2,-26) {\textbf{Place lock}};
% feed: sleeping, with request-reply slivers
\fill[navyL] (0,-2) rectangle (10*\s,2);
\foreach \t/\l in {1.88/{reply adds order 1},4.88/{reply adds order 2},7.88/{reply}} {
  \fill[navyM] (\t*\s-0.4,-2) rectangle (\t*\s+0.4,2);
  \node[f,anchor=south,text=navy] at (\t*\s,2.3) {\l};}
\foreach \t in {0.4,3.38,6.38} {\node[f,text=navy] at (\t*\s+5,0) {\code{sleep(3)}};}
\node[f,text=navy] at (8.88*\s+5,0) {\code{sleep(3)}};
% R lane
\fill[navy] (0,-14) rectangle (0.14*\s+0.6,-10);
\fill[navyL] (0.14*\s+0.6,-14) rectangle (2.14*\s,-10);
\fill[navy] (2.14*\s,-14) rectangle (2.25*\s+0.6,-10);
\fill[navyL] (2.25*\s+0.6,-14) rectangle (5.25*\s,-10);
\foreach \t in {0.64,1.14,1.64,2.75,3.25,3.75,4.25,4.75} {\draw[white,line width=0.7pt] (\t*\s,-14) -- (\t*\s,-10);}
\node[f,text=navy,anchor=west] at (5.25*\s+1.5,-12) {\code{rebalance()} returns both order IDs};
% found markers and causal arrows
\foreach \t/\r in {2.14/1.88,5.25/4.88} {
  \fill[navy] (\t*\s,-9.6) -- ++(-1.1,1.8) -- ++(2.2,0) -- cycle;
  \draw[ev] (\r*\s,-2.2) -- (\t*\s,-7.6);}
\node[f,text=navy,anchor=west] at (2.14*\s+1.5,-6.6) {order 1 found (2.0\,s)};
\node[f,text=navy,anchor=west] at (5.25*\s+1.5,-6.6) {order 2 found (3.0\,s)};
% annotations below R
\draw[decorate,decoration={brace,amplitude=1.2mm,mirror},greyD,line width=0.5pt] (0.14*\s,-15) -- (1.14*\s,-15);
\node[note,anchor=north west,align=left] at (0,-16.8) {order 1 sent;\\gate shut 1\,s};
\node[note,anchor=north west,align=left] at (2.14*\s-0.5,-15.4) {order 2 sent:\\gate open (2.1\,s)};
\node[note,anchor=north west,align=left] at (5.25*\s+1.5,-15.4) {white ticks: R looks in \code{hub.state} every 0.5\,s for up to 15\,s,\\then in REST \code{Orders}; if the order is in neither, R raises\\\code{AmbiguousState}; it never resends the order};
% place lock
\fill[navy] (0.07*\s,-28) rectangle (2.14*\s,-24);
\fill[navy] (2.14*\s,-28) rectangle (5.25*\s,-24);
\draw[white,line width=0.7pt] (2.14*\s,-28) -- (2.14*\s,-24);
\node[f,text=white] at ({(0.07+2.14)/2*\s},-26) {order 1};
\node[f,text=white] at ({(2.14+5.25)/2*\s},-26) {order 2: check, write, find};
% axis
\draw[greyD,line width=0.4pt] (0,-31) -- (10*\s,-31);
\foreach \t in {0,...,10} {\draw[greyD,line width=0.4pt] (\t*\s,-31) -- (\t*\s,-32); \node[f,text=greyD,anchor=north] at (\t*\s,-32) {\t};}
\node[f,text=greyD,anchor=north] at (5*\s,-35.2) {s after R starts};
\end{scope}
% ====================================================== legend
\begin{scope}[yshift=-120mm]
\foreach \c/\x in {navy/0,navyM/34,greyM/64,navyL/96} {\fill[\c] (\x,0) rectangle ++(3.5,2.4);}
\node[f,anchor=west] at (4.3,1.2) {Running on the thread};
\node[f,anchor=west] at (38.3,1.2) {Waiting for a reply};
\node[f,anchor=west] at (68.3,1.2) {Waiting for a lock};
\node[f,anchor=west] at (100.3,1.2) {Sleeping or looking in the state};
\fill[navy] (1.75,-2.6) -- ++(-1.1,1.8) -- ++(2.2,0) -- cycle; \node[f,anchor=west] at (4.3,-1.7) {order ID found};
\draw[ev] (34,-1.7) -- (37.5,-1.7); \node[f,anchor=west] at (38.3,-1.7) {feed reply changes \code{hub.state}};
\fill[navy] (96,-2.9) rectangle ++(3.5,2.4); \node[f,anchor=west] at (100.3,-1.7) {lock held (lock lanes)};
\end{scope}
\end{tikzpicture}}{The lock lanes use navy for ``held''. Measured: requests from Python take 31 to 39\ms{} (Section~\ref{sec:transport}); orders appeared in the feed 1 to 3\,s after \code{AddOrder}. Wrapper settings: feed request every 3\,s, look in the state every 0.5\,s, 1\,s between writes, 80 orders a day, renewal at the first call after 720\,s by the server's clock. Illustrative: the start times, the order of the coroutines and the feed reply times, chosen within the measured ranges. Gantt layout after Silberschatz et al.\ [7]. In panel (b), times in brackets run from sending an order to finding it, and from the end of the last write to the next write.}{DART Quant Team; code at \code{src/pytimefolio} (\code{auth}, \code{orders}, \code{hub})}


\lead{The code splits into 8 modules: 1 for each of the 5 layers of Section~\ref{sec:layers}, \code{state} for the feed's records, \code{client} to join them, and \code{errors} (\figref{modules}).} Each module imports only modules above it in the table, so the transport and authentication can be tested without a live feed, and the order arithmetic without a server. \code{errors} defines 1 exception per failure, all subclasses of \code{TimefolioError}: \code{AuthError} for a 401, \code{RequestError} for a 400, \code{ServerError} for a 5xx or a lost connection, \code{BlockedError} for a 403 or 429, \code{StaleFeed}, \code{RebalanceError} and \code{AmbiguousState}, the last for goal G5.

% OWNER: figs/modules.tex. Edit only this file.
\fig[wide]{modules}{Eight modules: one per protocol layer, plus the feed state, the client class and the exceptions}[Module layout of \code{pytimefolio}, in import order]{%
{\tabfont\begin{tabularx}{\linewidth}{>{\RaggedRight\arraybackslash}p{18mm} >{\RaggedRight\arraybackslash}X >{\RaggedRight\arraybackslash}p{24mm}}
\thd{Module} & \thd{Responsibility} & \thd{Rests on}\\
\code{errors} & \code{TimefolioError} and its subclasses: \code{AuthError} (401), \code{RequestError} (400), \code{ServerError} (5xx or a lost connection, with a flag for whether the request was sent), \code{BlockedError} (403 or 429), \code{StaleFeed}, \code{RebalanceError} (carries the IDs already placed), \code{AmbiguousState} & \ref{sec:design}\\\rowsep
\code{transport} & One \code{httpx.AsyncClient(http2=True)} reused for every REST request; \code{warm_up()}; server clock offset from the replies' \code{Date} headers; status codes turned into exceptions, with the first 403 or 429 on any request recorded; a GET repeated once after a connection or read/write failure (not after a timeout), a POST never & \ref{sec:transport}, \ref{sec:auth}\\\rowsep
\code{auth} & Renewal from the stored cookie jar at start, login only if it fails; lazy renewal at the first call past the renewal age or on a 401, one at a time under a lock with a re-check inside; login when the cookie jar is near expiry; cookie jar written atomically, readable by the owner only; password taken from the caller and never stored; Bearer header on every call & \ref{sec:auth}\\\rowsep
\code{rest} & \code{currBizDay} and the check against the KST date; portfolio rule; \code{Session}, \code{Universe}, \code{Trades}, \code{Violations}, \code{Orders} (\code{d0}, \code{d1}), \code{ValidateProduct} & \ref{sec:queries}, \ref{sec:orders}\\\rowsep
\code{state} & Snapshot replacement; delta merge keyed by \code{int(Id)}, including \code{bs_updated} into \code{bs}; null replies; index by \code{prodId}; other top-level fields (\code{cmpl}, \code{tags}) kept for the guidelines & \ref{sec:feed}, \ref{sec:guidelines}\\\rowsep
\code{hub} & SignalR negotiate, WebSocket with the token in a header, handshake, one request in flight with one 20\,s deadline that pings do not extend, pong to pings; time of the last reply, \code{StaleFeed} when it is older than 26\,s; reconnect after any failure with exponential backoff and full jitter, except a 403 or 429, which ends the feed & \ref{sec:feed}\\\rowsep
\code{orders} & Order body with the size quantised once; water-filling, sector check, target-to-order conversion net of working orders; order gate with the daily cap (a count file per business day), spacing between writes and a halt after any 403 or 429; pre-checks (fresh feed, business day, \code{ValidateProduct}, \code{weiLimit}); one lock across write and find; IDs already returned skipped and \code{AmbiguousState} when more than one new record matches; confirmation in the feed, then in REST \code{Orders}, never a resend of an order the server may have received; stop through \code{CancelOrderBatch} skipping finished orders; outcome when \code{Done} & \ref{sec:orders}, \ref{sec:oneshare}, \ref{sec:guidelines}\\\rowsep
\code{client} & \code{Timefolio}, an async context manager: start-up (session, portfolio, feed task, \code{Universe}, today's order count), state views, \code{buy}, \code{sell}, \code{stop}, \code{wait}, \code{nav()}, \code{share_weight()}, \code{plan()}, \code{rebalance()}, \code{guidelines()}, \code{orders_left()} & \ref{sec:design}, \ref{sec:live}\\
\end{tabularx}}}{Each module imports only modules above it. The other numeric parameters (intervals, timeouts, caps, renewal age) are listed in Appendix~\ref{app:decisions}. Water-filling: when a target breaks a stock's \code{weiLimit}, the excess is spread over the other targets in proportion to their weights, so the total is unchanged.}{DART Quant Team; code at \code{src/pytimefolio}}


\par\Needspace{6\baselineskip}
\claim{Operating the wrapper: each failure raises a named exception, and the logs hold no personal data}

\lead{Five signals tell the operator when to step in: 3 exceptions, the daily count and 1 guideline flag.} \code{StaleFeed} is raised by state reads, order placement and \code{wait()} when the feed's last reply is older than 26\,s; the feed task keeps reconnecting with full-jitter backoff (Section~\ref{sec:feed}), so a caller that catches it can retry once the feed answers again. \code{guidelines()["ok"]} is True only with 0 current server violations and at least 5\% turnover this week (Section~\ref{sec:guidelines}), so a caller can check one flag before and after a rebalance.

\lead{A 403 or 429 raises \code{BlockedError} and stops all writes for the rest of the process.} The organiser may block an IP address for bursts of traffic, and a 403 or 429 is the only sign of a block the wrapper can see. Reads still work, so the caller can record the state; orders resume only after a restart, which gives a person the chance to ask the organiser first.

\lead{An order the wrapper cannot confirm raises \code{AmbiguousState} and is never resent (Section~\ref{sec:orders}); a failed \code{rebalance()} raises \code{RebalanceError}.} \code{rebalance()} stops at the first failed order, and the exception's \code{placed} attribute holds the order IDs already sent, so the caller can wait on them or stop them.

\lead{The daily count stays at or under 80 and is synchronised with the server.} Each \code{AddOrder} adds 1 and each stopped order adds 1, since \code{CancelOrderBatch} counts every ID it stops; a request that never reached the server adds nothing. The count lives in 1 file per business day under \code{~/.config/pytimefolio}, written atomically under a file lock, and start-up and every \code{rebalance()} raise it to the larger of the local count and the server's record of the day's new orders and stops, so orders placed in the browser are counted too. By 10:34 on 24 Nov the counter showed 16 of 80 used (9 new orders, 7 stops). A request that would pass 80 raises before it is sent.


\lead{The logs record what the wrapper did and leave out every credential and personal field.} The wrapper writes through Python's \code{logging} module under the logger \code{pytimefolio} and its children \code{auth}, \code{hub} and \code{orders}; the caller chooses where the lines go. The example in the README sends them to standard error and raises \code{httpx} to WARNING, since \code{httpx} logs every URL at INFO. The wrapper logs the start-up line (portfolio, business day, orders used of 80), each login with the cookie expiry time, each renewal with the business day, each feed connection and failure with the backoff wait, each placed order (order ID, ticker, \code{wei}, \code{state}, \code{qty2ord}) and any block signal. It never logs the login reply, which carries the user's phone number, date of birth, \code{userNm} and \code{userNick}, nor access tokens, cookies or the password; the token travels in a header, so the URLs \code{httpx} logs carry none.

\lead{The cookie jar is the only credential the wrapper writes to disk, and logging out does not revoke it (\figref{security}).} The password is passed to the constructor, for example from environment variables, and never stored. The jar goes to \code{~/.config/pytimefolio/session.json} with mode 0600, written atomically at login and after each renewal; anyone who copies the file holds a working session for up to 7 days (Section~\ref{sec:auth}).

% OWNER: figs/security.tex. Edit only this file.
\fig[wide]{security}{The cookie jar is the one secret the wrapper stores}[Secrets and personal data that pass through \code{pytimefolio}, and how each is kept]{%
{\tabfont\begin{tabularx}{\linewidth}{>{\RaggedRight\arraybackslash}p{30mm} >{\RaggedRight\arraybackslash}p{62mm} >{\RaggedRight\arraybackslash}X}
\thd{Item} & \thd{What it allows or contains} & \thd{How the wrapper keeps it}\\
\hl Cookie jar & Renews access tokens for 7 days from login; it keeps renewing after \code{Logout}; its behaviour after a password change is untested & Stored in \code{~/.config/pytimefolio/session.json}, mode 0600, outside the repository; written atomically at login and after each renewal\\\rowsep
Password & Logs in and creates a new 7-day cookie jar & Not stored by the wrapper; the caller passes it to the constructor, for example from an environment variable\\\rowsep
Access token & Authenticates every REST and feed request for 900\,s & Kept in memory and sent only in the \code{Authorization} header; kept out of URLs and logs\\\rowsep
Login reply & Personal fields (contact details) & Never logged\\\rowsep
Position records & The account holder's name and nickname (\code{userNm}, \code{userNick}) & Never logged\\\rowsep
Repository & Public on GitHub & \code{.env}, \code{data/} and \code{logs/} are in \code{.gitignore}\\
\end{tabularx}}}{Shaded: the only secret written to disk. Token and cookie lifetimes are from Section~\ref{sec:auth}. The password-change test is listed with the open questions in Section~\ref{sec:validity}.}{DART Quant Team; code at \code{src/pytimefolio}}

